\chapter*{前言：先让机器完成一件小事}
\addcontentsline{toc}{chapter}{前言：先让机器完成一件小事}

操作系统教材很容易从一张完整架构图开始：进程、页表、调度器、驱动、文件系统和网络协议
同时出现在读者面前，每个名词都有定义，却没有一个机制显得非有不可。等到正文进入实现，
读者又会遇到另一种困难：启动日志、结构体字段、版本里程碑和测试结果接连出现，知道项目
“做到了什么”，仍然难以重建机器“为什么必须这样变化”。

本书选择另一条起点。我们先要求一台没有操作系统的机器完成一件很小的事：读出四个整数，
求和，把结果写回，再把控制权交还给一个已知位置。任务本身并不难，困难在于每一步都必须
落到可检查的机器事实上。第一条指令从哪里来？一串字节在什么时候可以被当作程序？处理器
等待存储响应时要保存什么？结果已经写入与整次运行已经结束是不是同一件事？只要这些问题
还有一个说不清楚，就不能用“程序运行了”把中间过程省略过去。

这项小任务同时规定了本书的解释方法。每一章开始时，先交代上一章留下的机器：哪些对象
已经存在，谁拥有它们，它们能沿什么路径改变。随后给出一个确实会发生的操作或失败，让
现有机器暴露一个边界。正文只增加解决该问题所需的最小机制，再沿数据流和控制流走到可见
结果；正常路径、部分完成、拒绝、回滚和恢复都在同一组状态中解释。章节结束时，机器已经
不同于开头，但仍只前进了一步。

贯穿全书的现实案例是一套自研 x86-64 系统。它从处理器复位入口开始，由自研 ROM 接管
机器，经磁盘装入启动阶段，建立长模式并把 ELF64 内核交给新的执行环境；内核随后接管
描述符、异常、内存、设备和用户边界，最终让多个用户程序通过描述符、管道、持久文件和
交互式 Shell 组合工作。案例还在继续演进：现有系统已经能完成一次完整的交互运行，但
固定资源、进程与线程身份、虚拟内存意图、命名空间和持久化恢复仍需要更清楚的对象模型。

这些能力不会在开篇被压成一张“大而全”的流程图。读者会先在小机器上看清一个边界，再
回到 x86-64 案例中辨认同一边界的真实形态。例如，教学机器用固定复位地址说明“稳定起点”
的含义；案例机器则必须从 x86 复位向量取得第一条指令。教学机器用一条四字节装入指令
说明请求何时提交；案例机器会把同一问题扩展到 ATA 事务、ELF 段装入和页表切换。前者
负责把因果关系缩小到可以手算，后者负责证明这些因果关系经得起真实体系结构的约束。

因此，本书既不是项目发布记录，也不是脱离实现的概念词典。项目提供纵向主线，已有的
系统知识提供横向深入；二者在一条可重放的运行轨迹中会合。读完一章，读者应当能够合上
书，从输入和初始状态开始，独立画出新增对象、关键状态转换、提交点和失败出口。如果只能
复述名词，却不能说明一次具体运行怎样穿过这些边界，就还没有真正掌握这一章。

\begin{statebox}
\textbf{本册样稿的范围。} 这里先验证叙述方法本身：从人工准备一项任务开始，逐步得到
自动装入、常驻协调和处理器强制返回。地址空间、用户权限、文件和网络尚未由当前问题
要求，因此不会提前塞进第一项任务。它们会在前一台机器确实失败时再出现。
\end{statebox}
